\chapter{Phrases at Japanese Convenience Store}

% ================= TITLE: CONVERSATION =================
\vspace{1em}
\begin{tcolorbox}[colback=boxframe, colframe=boxframe, rounded corners, boxrule=0pt, arc=3mm, left=2mm, top=2mm, bottom=2mm]
    \Large \bfseries \color{white} 🏪 1. The Ultimate Konbini Flow (Alur Kasir)
\end{tcolorbox}
\vspace{1em}

% ================= PAIR 1 =================
\begin{convobox}{1}
    \noindent
    \jpword{つぎ}{次}{Tsugi} \jpkana{の}{no} 
    \jpword{かた}{方}{kata} \jpkana{どうぞ。}{douzo.}

    \vspace{1.5em}
    \noindent{\Large \color{articolor} \textbf{Arti:} Orang berikutnya, silakan (maju).}
\end{convobox}

\begin{convobox}{2}
    \noindent
    \jpkana{いらっしゃいませ。}{Irasshaimase.} 
    \jpkana{お}{o}\jpword{あず}{預}{azu}\jpkana{かりいたします。}{kari itashimasu.}

    \vspace{1.5em}
    \noindent{\Large \color{articolor} \textbf{Arti:} Selamat datang. Saya akan menerima/menyimpan (barang belanjaan) Anda.}
\end{convobox}

\begin{lessonbox}{⛩️ Panduan Turis: Garis Antrean \& Panggilan Kasir!}
    Halo Ayah, Bunda, dan teman-teman! Di Jepang, **kita tidak boleh langsung maju dan meletakkan barang di meja kasir**. Kamu harus berdiri di belakang garis atau stiker jejak kaki di lantai. \\
    
    Tunggu sampai kasir menatapmu dan berkata \textbf{「次の方どうぞ」 (Tsugi no kata douzo - Orang berikutnya silakan)}. Setelah mendengar kata ajaib itu, barulah kamu boleh maju ke meja kasir!
\end{lessonbox}

% ================= PAIR 2 =================
\begin{convobox}{3}
    \noindent
    \jpkana{お}{O}\jpword{べんとう}{弁当}{bentou} \jpword{あたた}{温}{atata}\jpkana{めますか？}{memasu ka?}

    \vspace{1.5em}
    \noindent{\Large \color{articolor} \textbf{Arti:} Apakah kotak bekalnya (Bento) mau dipanaskan?}
\end{convobox}

\begin{convobox}{4}
    \noindent
    \jpkana{はい、}{Hai,} \jpkana{お}{o}\jpword{ねが}{願}{nega}\jpkana{いします。}{ishimasu.}

    \vspace{1.5em}
    \noindent{\Large \color{articolor} \textbf{Arti:} Ya, tolong.}
\end{convobox}

\begin{lessonbox}{🍱 Panduan Turis: Memanaskan Makanan}
    Minimarket Jepang punya makanan instan terenak di dunia! Kalau kamu membeli Nasi (Bento), Pasta, atau Onigiri tertentu, kasir otomatis akan bertanya \textit{"Atatamemasu ka?"} (Mau dipanaskan?). \\
    
    Jawab dengan lantang: \textbf{「はい、お願いします」(Hai, onegaishimasu)} dan mereka akan memasukkannya ke microwave khusus berkecepatan tinggi!
\end{lessonbox}

% ================= PAIR 3 =================
\begin{convobox}{5}
    \noindent
    \jpkana{お}{O}\jpword{はし}{箸}{hashi} \jpkana{ご}{go}\jpword{りよう}{利用}{riyou} \jpkana{ですか？}{desu ka?}

    \vspace{1.5em}
    \noindent{\Large \color{articolor} \textbf{Arti:} Apakah Anda menggunakan (membutuhkan) sumpit?}
\end{convobox}

\begin{convobox}{6}
    \noindent
    \jpkana{はい、}{Hai,} \jpkana{お}{o}\jpword{ねが}{願}{nega}\jpkana{いします。}{ishimasu.}

    \vspace{1.5em}
    \noindent{\Large \color{articolor} \textbf{Arti:} Ya, tolong.}
\end{convobox}

\begin{lessonbox}{🥢 Panduan Turis: Alat Makan Sesuai Kebutuhan}
    Kadang kasir tidak menyebut "O-hashi" (Sumpit). Kalau kamu beli yogurt, mereka akan bilang \textbf{「スプーン」 (Supuun - Sendok)}. Kalau beli pasta, mereka bilang \textbf{「フォーク」 (Fooku - Garpu)}. Dengarkan awalannya ya, dan cukup jawab "Hai, onegaishimasu" untuk mendapatkannya gratis!
\end{lessonbox}

% ================= PAIR 4 =================
\begin{convobox}{7}
    \noindent
    \jpkana{ポイントカード}{Pointo kaado} \jpkana{お}{o}\jpword{も}{持}{mo}\jpkana{ちですか？}{chi desu ka?}

    \vspace{1.5em}
    \noindent{\Large \color{articolor} \textbf{Arti:} Apakah Anda membawa kartu poin?}
\end{convobox}

\begin{convobox}{8}
    \noindent
    \jpkana{いえ、}{Ie,} \jpword{も}{持}{mo}\jpkana{っていません。}{tte imasen.}

    \vspace{1.5em}
    \noindent{\Large \color{articolor} \textbf{Arti:} Tidak, saya tidak membawanya.}
\end{convobox}

\begin{lessonbox}{💳 Panduan Turis: Lewati Saja Bagian Ini!}
    Turis tidak butuh kartu poin (T-Point, d-Point, Ponta). Agar prosesnya cepat, kamu bisa memotong dengan sopan dengan menggelengkan kepala dan tersenyum sambil bilang \textbf{「いえ、大丈夫です」 (Ie, daijoubu desu - Tidak usah/Tidak ada)}.
\end{lessonbox}

% ================= PAIR 5 =================
\begin{convobox}{9}
    \noindent
    \jpkana{そちらの}{Sochira no} 
    \jpword{ねんれい}{年齢}{nenrei} \jpword{かくにん}{確認}{kakunin} \jpkana{ボタンを}{botan o} 
    \jpword{お}{押}{o}\jpkana{していただけますでしょうか？}{shite itadakemasu deshou ka?}

    \vspace{1.5em}
    \noindent{\Large \color{articolor} \textbf{Arti:} Bisakah Anda menekan tombol konfirmasi umur yang ada di sebelah sana?}
\end{convobox}

\begin{lessonbox}{🚨 PANDUAN KRITIS: Jebakan Layar Sentuh!}
    \textbf{Perhatian khusus untuk orang tua!} Jika Ayah/Bunda membeli bir, minuman alkohol kalengan (seperti Strong Zero, Chuhai), atau rokok, tiba-tiba mesin kasir akan berbunyi dan kasir menunjuk layar sentuh ke arahmu. \\
    
    Itu adalah **Nenrei Kakunin (Konfirmasi Umur)**. Di Jepang, kamu harus berusia di atas 20 tahun. Di layar akan muncul tulisan \textbf{「私は20歳以上です」 (Saya di atas 20 tahun)}. \textbf{KAMU HARUS MENYENTUH TOMBOL DI LAYAR ITU!} Jika tidak disentuh, kasir tidak akan bisa melanjutkan transaksi. Sentuh saja layarnya, dan transaksi lanjut!
\end{lessonbox}

% ================= PAIR 6 =================
\begin{convobox}{10}
    \noindent
    \jpword{ふくろ}{袋}{Fukuro} \jpkana{ご}{go}\jpword{りよう}{利用}{riyou} \jpkana{ですか？}{desu ka?}

    \vspace{1.5em}
    \noindent{\Large \color{articolor} \textbf{Arti:} Apakah Anda menggunakan kantong plastik?}
\end{convobox}

\begin{convobox}{11}
    \noindent
    \jpkana{はい、}{Hai,} \jpkana{お}{o}\jpword{ねが}{願}{nega}\jpkana{いします。}{ishimasu.}

    \vspace{1.5em}
    \noindent{\Large \color{articolor} \textbf{Arti:} Ya, tolong.}
\end{convobox}

\begin{lessonbox}{🛍️ Panduan Turis: Kantong Plastik (Fukuro)}
    Selalu ingat kata kunci \textbf{FUKURO}. Karena tanganmu mungkin sudah penuh memegang payung atau kamera, meminta plastik adalah ide bagus meskipun berbayar.
\end{lessonbox}

% ================= PAIR 7 =================
\begin{convobox}{12}
    \noindent
    \jpword{いちまい}{1枚}{Ichi-mai} \jpword{さんえん}{3円}{san-en} 
    \jpkana{ですが}{desu ga} 
    \jpkana{よろしいでしょうか？}{yoroshii deshou ka?}

    \vspace{1.5em}
    \noindent{\Large \color{articolor} \textbf{Arti:} 1 lembarnya seharga 3 yen, apakah tidak apa-apa (apakah Anda setuju)?}
\end{convobox}

\begin{convobox}{13}
    \noindent
    \jpkana{はい。}{Hai.}

    \vspace{1.5em}
    \noindent{\Large \color{articolor} \textbf{Arti:} Ya.}
\end{convobox}

\begin{lessonbox}{💴 Panduan Turis: Kasir yang Transparan}
    Kasir Jepang sangat jujur. Jika ada biaya tambahan (seperti plastik 3 yen sampai 5 yen), mereka akan selalu meminta persetujuanmu dengan kata \textbf{「よろしいでしょうか」 (Yoroshii deshou ka - Apakah tidak apa-apa?)}. Jawab saja "Hai!" dengan percaya diri.
\end{lessonbox}

% ================= PAIR 8 =================
\begin{convobox}{14}
    \noindent
    \jpkana{お}{O}\jpword{かいけい}{会計}{kaikei} 
    \jpword{ななひゃくななじゅうさんえん}{773円}{nana-hyaku nana-juu san-en} 
    \jpkana{でございます。}{de gozaimasu.}

    \vspace{1.5em}
    \noindent{\Large \color{articolor} \textbf{Arti:} Total pembayarannya adalah 773 yen.}
\end{convobox}

\vspace{2em}
\begin{tcolorbox}[colback=boxframe, colframe=boxframe, rounded corners, boxrule=0pt, arc=3mm, left=2mm, top=2mm, bottom=2mm]
    \Large \bfseries \color{white} 💬 2. Alternatif Jawaban (Saat Berubah Pikiran)
\end{tcolorbox}
\vspace{1em}

\begin{convobox}{15}
    \noindent
    \jpkana{いえ、}{Ie,} \jpword{だいじょうぶ}{大丈夫}{daijoubu} \jpkana{です。}{desu.}

    \vspace{1.5em}
    \noindent{\Large \color{articolor} \textbf{Arti:} Tidak, tidak usah (Tidak apa-apa).}
\end{convobox}

\begin{convobox}{16}
    \noindent
    \jpkana{やっぱり}{Yappari} 
    \jpword{い}{要}{i}\jpkana{りません。}{rimasen.}

    \vspace{1.5em}
    \noindent{\Large \color{articolor} \textbf{Arti:} Sepertinya aku tidak jadi butuh.}
\end{convobox}

\begin{lessonbox}{🔄 Panduan Turis: "Eh, Gak Jadi Deh!"}
    Bagaimana kalau kamu awalnya minta plastik, lalu ingat kamu bawa tas kain sendiri? 
    Gunakan kata ajaib \textbf{やっぱり (Yappari)} yang artinya "Sebenarnya/Gak jadi deh". Lalu sambung dengan \textbf{要りません (Irimasen - Saya tidak butuh)}. 
    
    "Yappari irimasen!" = "Eh, gak jadi butuh deh plastikannya!". Kasir akan langsung paham dan membatalkan plastiknya!
\end{lessonbox}


% --- SUMMARY BOX ---
\vspace{2em}
\begin{tcolorbox}[
    colback=blue!5!white,
    colframe=blue!80!black,
    boxrule=3pt,
    arc=5mm,
    title={\huge\bfseries \sffamily 🌟 Rangkuman Kelulusan Panduan Jepang! 🌟},
    fonttitle=\color{white},
    attach boxed title to top center={yshift=-4mm},
    boxed title style={colback=blue!80!black, rounded corners},
    left=5mm, right=5mm, top=7mm, bottom=5mm
]
\Large \textbf{SELAMAT! Kamu Sudah Lulus dan Siap Berlibur ke Jepang!}
\begin{itemize}
    \item \textbf{Antre dengan Benar:} Tunggu aba-aba \textit{Tsugi no kata douzo} sebelum maju ke meja kasir.
    \item \textbf{Siap Dipanaskan:} Bento/Nasi akan ditanya \textit{Atatamemasu ka?}. Jawab \textit{Hai, onegaishimasu!}
    \item \textbf{Jangan Lupa Tombol Layar:} Membeli alkohol = Sentuh layar \textit{Nenrei Kakunin} (Konfirmasi Umur 20+).
    \item \textbf{Ubah Pikiran:} Gunakan \textit{Yappari Irimasen} kalau kamu batal meminta kantong plastik.
\end{itemize}

Terima kasih sudah belajar dari Bab 1 hingga Bab 12 bersama buku **Andi - Japanese Handbook**. Perjalananmu mempelajari bahasa dan budaya Jepang sangat luar biasa. Bawalah buku ini bersamamu saat ke Jepang nanti, dan jangan ragu untuk berbicara dengan penduduk lokal! \textbf{気をつけて、行ってらっしゃい！ (Ki o tsukete, Itterasshai - Hati-hati di jalan dan selamat bersenang-senang!)}
\end{tcolorbox}

\newpage